\documentclass[10pt,a4paper]{amsart}
\usepackage{amsmath,amssymb,amsthm,graphicx}
% Retained preamble for the Stacks Project chapter {CH} (shared project preamble).
% The chapter file pinned-original-source.tex only carries \input{{preamble}} on its first line,
% so no preamble text lives inside the chapter itself.  The Stacks Project uses one single shared
% preamble for every chapter: the snapshot stages it three times (sources/stacks-chow/preamble.tex,
% sources/stacks-etale-cohomology/preamble.tex and focused-seeds/AGP/sources/stacks-snapshot-20261011/
% preamble.tex) and those copies are identical in size.  The block below is copied mechanically,
% line for line, from sources/stacks-chow/preamble.tex lines 14-18 and 154-217: the Xy-pic setup and
% the native theorem declarations together with the mathematical macros.  Nothing was added, removed
% or redefined; the wrapper class and the amsmath/amssymb/amsthm/graphicx base packages are supplied
% by the assembler.
% For commutative diagrams we use Xy-pic
% Chapter: stacks-spaces-more-morphisms.tex -- More on Morphisms of Spaces.
\usepackage[all]{xy}

% We use 2cell for 2-commutative diagrams.
\xyoption{2cell}

% Theorem environments.
%
\theoremstyle{plain}
\newtheorem{theorem}[subsection]{Theorem}
\newtheorem{proposition}[subsection]{Proposition}
\newtheorem{lemma}[subsection]{Lemma}

\theoremstyle{definition}
\newtheorem{definition}[subsection]{Definition}
\newtheorem{example}[subsection]{Example}
\newtheorem{exercise}[subsection]{Exercise}
\newtheorem{situation}[subsection]{Situation}

\theoremstyle{remark}
\newtheorem{remark}[subsection]{Remark}
\newtheorem{remarks}[subsection]{Remarks}

\numberwithin{equation}{subsection}

% Macros
%
\def\lim{\mathop{\mathrm{lim}}\nolimits}
\def\colim{\mathop{\mathrm{colim}}\nolimits}
\def\Spec{\mathop{\mathrm{Spec}}}
\def\Hom{\mathop{\mathrm{Hom}}\nolimits}
\def\Ext{\mathop{\mathrm{Ext}}\nolimits}
\def\SheafHom{\mathop{\mathcal{H}\!\mathit{om}}\nolimits}
\def\SheafExt{\mathop{\mathcal{E}\!\mathit{xt}}\nolimits}
\def\Sch{\mathit{Sch}}
\def\Mor{\mathop{\mathrm{Mor}}\nolimits}
\def\Ob{\mathop{\mathrm{Ob}}\nolimits}
\def\Sh{\mathop{\mathit{Sh}}\nolimits}
\def\NL{\mathop{N\!L}\nolimits}
\def\CH{\mathop{\mathrm{CH}}\nolimits}
\def\proetale{{pro\text{-}\acute{e}tale}}
\def\etale{{\acute{e}tale}}
\def\QCoh{\mathit{QCoh}}
\def\Ker{\mathop{\mathrm{Ker}}}
\def\Im{\mathop{\mathrm{Im}}}
\def\Coker{\mathop{\mathrm{Coker}}}
\def\Coim{\mathop{\mathrm{Coim}}}

% Boxtimes
%
\DeclareMathSymbol{\boxtimes}{\mathbin}{AMSa}{"02}

%
% Macros for moduli stacks/spaces
%
\def\QCohstack{\mathcal{QC}\!\mathit{oh}}
\def\Cohstack{\mathcal{C}\!\mathit{oh}}
\def\Spacesstack{\mathcal{S}\!\mathit{paces}}
\def\Quotfunctor{\mathrm{Quot}}
\def\Hilbfunctor{\mathrm{Hilb}}
\def\Curvesstack{\mathcal{C}\!\mathit{urves}}
\def\Polarizedstack{\mathcal{P}\!\mathit{olarized}}
\def\Complexesstack{\mathcal{C}\!\mathit{omplexes}}
% \Pic is the operator that assigns to X its picard group, usage \Pic(X)
% \Picardstack_{X/B} denotes the Picard stack of X over B
% \Picardfunctor_{X/B} denotes the Picard functor of X over B
\def\Pic{\mathop{\mathrm{Pic}}\nolimits}
\def\Picardstack{\mathcal{P}\!\mathit{ic}}
\def\Picardfunctor{\mathrm{Pic}}
\def\Deformationcategory{\mathcal{D}\!\mathit{ef}}
\newenvironment{reference}{}{}
\newenvironment{slogan}{}{}
\newenvironment{history}{}{}
\usepackage{cleveref}
\crefname{originallemma}{original lemma}{original lemmas}
\Crefname{originallemma}{Original lemma}{Original lemmas}
\makeatletter
\newlabel{lemma-blow-up-resolution-property}{{0GV5}{source}{}{sourcepointer-1}{}}
\newlabel{lemma-blow-up-resolution-property@cref}{{[originallemma][0][]0GV5}{[1][0][]source}{}{}{}}
\makeatother
\crefname{originallemma}{original lemma}{original lemmas}
\Crefname{originallemma}{Original lemma}{Original lemmas}
\makeatletter
\newlabel{spaces-definition-base-change}{{03I5}{source}{}{sourcepointer-2}{}}
\newlabel{spaces-definition-base-change@cref}{{[originallemma][0][]03I5}{[1][0][]source}{}{}{}}
\makeatother
\crefname{originallemma}{original lemma}{original lemmas}
\Crefname{originallemma}{Original lemma}{Original lemmas}
\makeatletter
\newlabel{spaces-limits-proposition-approximate}{{07SU}{source}{}{sourcepointer-3}{}}
\newlabel{spaces-limits-proposition-approximate@cref}{{[originallemma][0][]07SU}{[1][0][]source}{}{}{}}
\makeatother
\crefname{originallemma}{original lemma}{original lemmas}
\Crefname{originallemma}{Original lemma}{Original lemmas}
\makeatletter
\newlabel{spaces-perfect-lemma-resolution-property-goes-up-affine}{{0GUU}{source}{}{sourcepointer-4}{}}
\newlabel{spaces-perfect-lemma-resolution-property-goes-up-affine@cref}{{[originallemma][0][]0GUU}{[1][0][]source}{}{}{}}
\makeatother
\crefname{originallemma}{original lemma}{original lemmas}
\Crefname{originallemma}{Original lemma}{Original lemmas}
\makeatletter
\newlabel{spaces-properties-definition-separated}{{03BS}{source}{}{sourcepointer-5}{}}
\newlabel{spaces-properties-definition-separated@cref}{{[originallemma][0][]03BS}{[1][0][]source}{}{}{}}
\makeatother
\title{Every qcqs algebraic space admits a finite filtration by closed subspaces whose strata become resolution spaces after admissible blowing ups}
\author{The Stacks Project Authors}
\date{stacks-project:spaces-more-morphisms.tex:local-snapshot-20261011; GNU FDL 1.2 or later}
\begin{document}
\maketitle
\noindent\textit{Editorial note.} Statement span is lines 13880-13892 of sources/stacks-spaces-more-morphisms/pinned-original-source.tex; proof span is lines 13894-13933, the complete author proof (38 body lines).\par\medskip
\noindent\textit{Editorial note.} Navigation stubs: 5 cited-but-not-included labels are supplied as non-load-bearing pointers with their pinned tags where the snapshot tag index has them.\par\medskip
\noindent\textit{Editorial note.} extra\_packages is empty on purpose: the retained preamble already loads xy with options and a second declaration clashes in the assembler wrapper.\par\medskip
\noindent\textit{Editorial note.} Text is retained verbatim, including the author's own typos and placeholders.\par\medskip
\noindent\textit{Editorial note.} The proof's first paragraph marks pullback details omitted; the induction core is written in full inside the retained span.\par\medskip
\noindent\textit{Allowed standard prerequisites (editorial).}\begin{itemize}
\item Explicit prerequisite boundary. Absolute Noetherian approximation for qcqs spaces (Limits of Spaces, proposition-approximate), ascent of the resolution property along affine morphisms (Derived Categories of Spaces, lemma-resolution-property-goes-up-affine), and the single-step blowup lemma of the same chapter (lemma-blow-up-resolution-property).
\end{itemize}\medskip
\ifdefined\hypertarget\hypertarget{sourcepointer-1}{}\fi
\noindent\textit{Source reference.} Original lemma 0GV5: sources/stacks-spaces-more-morphisms/pinned-original-source.tex, source lines 13833--13839. Cited by the retained statement or proof but not part of the collected spans; supplied only as a navigation pointer. The argument does not depend on its content.\par\medskip
\ifdefined\hypertarget\hypertarget{sourcepointer-2}{}\fi
\noindent\textit{Source reference.} Original lemma 03I5: sources/stacks-spaces-more-morphisms/pinned-original-source.tex, source lines 13933--13933. Cited by the retained statement or proof but not part of the collected spans; supplied only as a navigation pointer. The argument does not depend on its content.\par\medskip
\ifdefined\hypertarget\hypertarget{sourcepointer-3}{}\fi
\noindent\textit{Source reference.} Original lemma 07SU: sources/stacks-spaces-more-morphisms/pinned-original-source.tex, source lines 13933--13933. Cited by the retained statement or proof but not part of the collected spans; supplied only as a navigation pointer. The argument does not depend on its content.\par\medskip
\ifdefined\hypertarget\hypertarget{sourcepointer-4}{}\fi
\noindent\textit{Source reference.} Original lemma 0GUU: sources/stacks-spaces-more-morphisms/pinned-original-source.tex, source lines 13933--13933. Cited by the retained statement or proof but not part of the collected spans; supplied only as a navigation pointer. The argument does not depend on its content.\par\medskip
\ifdefined\hypertarget\hypertarget{sourcepointer-5}{}\fi
\noindent\textit{Source reference.} Original lemma 03BS: sources/stacks-spaces-more-morphisms/pinned-original-source.tex, source lines 13933--13933. Cited by the retained statement or proof but not part of the collected spans; supplied only as a navigation pointer. The argument does not depend on its content.\par\medskip
\begin{lemma}
\label{lemma-blow-up-resolution-property-filtered}
Let $S$ be a scheme. Let $X$ be a quasi-compact and quasi-separated
algebraic space over $S$. There exists a $t \geq 0$ and closed
subspaces
$$
X \supset Z_0 \supset Z_1 \supset \ldots \supset Z_t = \emptyset
$$
such that $Z_i \to X$ is of finite presentation,
$Z_0 \subset X$ is a thickening, and for each $i = 0, \ldots t - 1$
there exists a $(Z_i \setminus Z_{i - 1})$-admissible
blowing up $Z'_i \to Z_i$ such that $Z'_i$ has the resolution property.
\end{lemma}

\begin{proof}
In this paragraph we use absolute Noetherian approximation to reduce to the
case of algebraic spaces of finite presentation over $\Spec(\mathbf{Z})$.
We may view $X$ as an algebraic space over $\Spec(\mathbf{Z})$, see
Spaces, Definition \ref{spaces-definition-base-change} and
Properties of Spaces, Definition \ref{spaces-properties-definition-separated}.
Thus we may apply Limits of Spaces, Proposition
\ref{spaces-limits-proposition-approximate}.
It follows that we can find an affine morphism $X \to X_0$
with $X_0$ of finite presentation over $\mathbf{Z}$.
If we can prove the lemma for $X_0$, then we can pull back
the stratification and the centers of the blowing ups to $X$
and get the result for $X$; this uses that the resolution property
goes up along affine morphisms (Derived Categories of Spaces,
Lemma \ref{spaces-perfect-lemma-resolution-property-goes-up-affine}) and
that the strict transform of an affine morphism is affine -- details omitted.
This reduces us to the case discussed in the next paragraph.

\medskip\noindent
Assume $X$ is of finite presentation over $\mathbf{Z}$.
Then $X$ is Noetherian and $|X|$ is a Noetherian topological
space (with finitely many irreducible components) of finite dimension.
Hence we may use induction on $\dim(|X|)$.
By Lemma \ref{lemma-blow-up-resolution-property}
there exists a dense open $U \subset X$ and a $U$-admissible
blowing up $X' \to X$ such that $X'$ has the resolution property.
Set $Z_0 = X$ and let $Z_1 \subset X$ be the reduced closed subspace
with $|Z_1| = |X| \setminus |U|$.
By induction we find an integer $t \geq 0$ and a filtration
$$
Z_1 \supset Z_{1, 0} \supset Z_{1, 1} \supset \ldots
\supset Z_{1, t} = \emptyset
$$
by closed subspaces, where $Z_{1, 0} \to Z_1$ is a thickening
and there exist $(Z_{1, i} \setminus Z_{1, i + 1})$-admissible
blowing ups $Z'_{1, i} \to Z_{1, i}$ such that $Z'_{1, i}$
has the resolution property. Since $Z_1$ is reduced, we have $Z_1 = Z_{1, 0}$.
Hence we can set $Z_i = Z_{1, i - 1}$ and $Z'_i = Z'_{1, i - 1}$
for $i \geq 1$ and the lemma is proved.
\end{proof}

\end{document}
